\documentclass[12pt,a4paper]{article}
\usepackage[utf8]{inputenc}
\usepackage[T1]{fontenc}
\usepackage{lmodern}
\usepackage{amsmath,amssymb,amsthm}
\usepackage[top=2.5cm,bottom=2.5cm,left=2.8cm,right=2.8cm]{geometry}
\usepackage{parskip}
\usepackage{booktabs}

\newtheorem{theorem}{Théorème}
\newtheorem{lemma}[theorem]{Lemme}
\newtheorem{corollary}[theorem]{Corollaire}
\newtheorem{proposition}[theorem]{Proposition}
\newtheorem{definition}[theorem]{Définition}
\theoremstyle{remark}
\newtheorem{remark}[theorem]{Remarque}
\newtheorem{hypothese}[theorem]{Hypothèse}

\newcommand{\Inc}{\ensuremath{\mathrm{Inc}}}
\newcommand{\Exc}{\ensuremath{\mathrm{Exc}}}
\newcommand{\vbar}{\bar{v}}

\begin{document}

\begin{center}
{\LARGE\bfseries Convergence d'une règle locale à quatre cas pour l'apprentissage de conjonctions}
\end{center}
\vspace{1em}

\section{Cadre et notations}

Soit $(X_t,Y_t)_{t\ge1}\overset{\mathrm{i.i.d.}}{\sim}P$ de loi jointe
$P(X,Y)$ arbitraire, avec $X=(X_1,\ldots,X_n)\in\{0,1\}^n$ et $Y\in\{0,1\}$.
Pour chaque variable $X_i$, deux automates à $2N$ états, $v_i$ et $\vbar_i$,
gouvernent respectivement le littéral $x_i$ et sa négation $\neg x_i$ :
\[
  \mathrm{Include}(s) = \mathbf 1[s>N], \qquad s\in\{1,\ldots,2N\}.
\]
La clause induite par ces automates est
\[
  C(x) = \bigwedge_{i:\,v_i>N} x_i \;\wedge\; \bigwedge_{i:\,\vbar_i>N} \neg x_i.
\]
Les mouvements élémentaires d'un automate sont réfléchis aux bornes :
\[
  s \xrightarrow{\ \Inc\ } \min(s+1,2N), \qquad
  s \xrightarrow{\ \Exc\ } \max(s-1,1).
\]

\begin{definition}[Règle à quatre cas]\label{def:regle}
Avec $a,b,c,d\in(0,1]$, chaque automate est mis à jour à partir de la seule
paire locale $(x_i,y)$ (jamais de la sortie de la clause ni de l'état des
autres automates) selon :
\begin{center}
\begin{tabular}{cccc}
\toprule
$x_i$ & $y$ & Action sur $v_i$ & Action sur $\vbar_i$ \\
\midrule
0 & 0 & vers $\Inc$ avec proba.\ $a$ & vers $\Exc$ avec proba.\ $c$ \\
0 & 1 & vers $\Exc$ avec proba.\ $b$ & vers $\Inc$ avec proba.\ $d$ \\
1 & 0 & vers $\Exc$ avec proba.\ $c$ & vers $\Inc$ avec proba.\ $a$ \\
1 & 1 & vers $\Inc$ avec proba.\ $d$ & vers $\Exc$ avec proba.\ $b$ \\
\bottomrule
\end{tabular}
\end{center}
\end{definition}

\section{Dynamique d'un automate isolé}

\begin{proposition}[Probabilités de transition]\label{prop:pv}
Pour toute loi $P$, en notant $A=P(X_i{=}0,Y{=}0)$, $B=P(X_i{=}1,Y{=}0)$,
$C=P(X_i{=}0,Y{=}1)$, $D=P(X_i{=}1,Y{=}1)$ :
\begin{align*}
  p_i^+ &= aA+dD, & p_i^- &= bC+cB; \\
  \bar p_i^+ &= dC+aB, & \bar p_i^- &= cA+bD.
\end{align*}
Pour $p_i^+,p_i^->0$, $(v_{i,t})_t$ est une chaîne de Markov naissance--mort
réfléchie sur $\{1,\ldots,2N\}$ (de même pour $\vbar_{i,t}$).
\end{proposition}

\begin{proof}
Somme des contributions de la Définition~\ref{def:regle} sur les quatre cas
$(x_i,y)$, pondérées par leur probabilité : $v_i$ est poussé vers $\Inc$ avec
probabilité $a$ sur l'évènement $(X_i{=}0,Y{=}0)$ (probabilité $A$) et avec
probabilité $d$ sur $(X_i{=}1,Y{=}1)$ (probabilité $D$), d'où $p_i^+=aA+dD$ ;
de même pour les trois autres.
\end{proof}

\section{Loi stationnaire exacte}

\begin{theorem}\label{thm:stat}
Soit $\rho=p^+/p^-$, avec $p^+,p^->0$ (indice $i$ omis). La chaîne $(v_t)$
est irréductible et apériodique sur $\{1,\ldots,2N\}$, donc
$\mathcal L(v_t)\xrightarrow[t\to\infty]{}\pi$, de loi stationnaire
\[
  \pi(k) = \frac{(1-\rho)\rho^{k-1}}{1-\rho^{2N}}\ (\rho\ne1), \qquad
  \pi(k) = \frac1{2N}\ (\rho=1), \qquad
  \pi(\Inc) = \frac{\rho^N}{1+\rho^N}.
\]
Si $\rho>1$ : $\pi(\Inc)\to1$ et $1-\pi(\Inc)\le\rho^{-N}$. Si $\rho<1$ :
$\pi(\Exc)\to1$ symétriquement. De plus $\frac1T\sum_{t\le T}\mathbf1[v_t>N]\to\pi(\Inc)$ p.s.
\end{theorem}

\begin{proof}
\emph{Irréductibilité.} Chaîne de naissance-mort à transitions intérieures
strictement positives.

\emph{Noyau explicite.} Pour $2\le k\le 2N-1$ : $P(k,k{+}1)=p^+$,
$P(k,k{-}1)=p^-$, $P(k,k)=1-p^+-p^-$. Aux bornes : $P(1,1)=1-p^+$,
$P(1,2)=p^+$, $P(2N,2N)=1-p^-$, $P(2N,2N{-}1)=p^-$.

\emph{Apériodicité.} $P(1,1)=1-p^+\ge p^->0$ : l'état $1$ a un cycle de
longueur $1$, donc une période $1$ ; par irréductibilité, toute la chaîne
est apériodique.

\emph{Forme géométrique.} Bilan détaillé $\pi(k)p^+=\pi(k{+}1)p^-$ donne
$\pi(k{+}1)/\pi(k)=\rho$ constant, soit $\pi(k)=\pi(1)\rho^{k-1}$ ;
normalisation $\sum_{k=1}^{2N}\pi(k)=1$ donne $\pi(1)=(1-\rho)/(1-\rho^{2N})$,
d'où les formules annoncées. $\pi(\Inc)=\sum_{k=N+1}^{2N}\pi(k)=\rho^N/(1+\rho^N)$.
La convergence p.s.\ de la moyenne temporelle est le théorème ergodique
standard pour les chaînes de Markov finies irréductibles apériodiques.
\end{proof}

\begin{corollary}[Cas absorbant]\label{cor:absorb}
Si $p^-=0,p^+>0$ : absorption p.s.\ en $2N$ en temps fini. Si $p^+=0,p^->0$ :
symétrique en $1$. Si $p^+=p^-=0$ : l'automate reste figé à son état initial.
\end{corollary}

\section{Régimes de sélection d'un littéral}

\begin{lemma}[Somme des dérives]\label{lem:sum}
Avec $\Delta=p^+-p^-$, $\bar\Delta=\bar p^+-\bar p^-$ :
\[
  \Delta+\bar\Delta = (a-c)\,P(Y{=}0) + (d-b)\,P(Y{=}1).
\]
\end{lemma}

\begin{proof}
\begin{align*}
\Delta+\bar\Delta &= (aA+dD-bC-cB) + (dC+aB-cA-bD) \\
&= (a-c)(A+B) + (d-b)(C+D) = (a-c)P(Y{=}0)+(d-b)P(Y{=}1),
\end{align*}
puisque $A+B=P(Y{=}0)$ et $C+D=P(Y{=}1)$.
\end{proof}

\begin{hypothese}\label{hyp:H}
$a\le c$ et $d\le b$.
\end{hypothese}

\begin{corollary}[Non-contradiction]\label{cor:excl}
Sous l'Hypothèse~\ref{hyp:H}, pour toute loi $P(X,Y)$ : $\Delta+\bar\Delta\le0$.
En particulier, $\Delta>0$ et $\bar\Delta>0$ ne peuvent avoir lieu
simultanément : si $\Delta>0$ alors $\bar\Delta\le-\Delta<0$.
\end{corollary}

\begin{proof}
$(a-c)\le0$ et $(d-b)\le0$ par l'Hypothèse~\ref{hyp:H}, et
$P(Y{=}0),P(Y{=}1)\ge0$, donc la somme du Lemme~\ref{lem:sum} est $\le0$. Si
$\Delta>0$, alors $\bar\Delta=(\Delta+\bar\Delta)-\Delta\le-\Delta<0$.
\end{proof}

\begin{theorem}[Trois régimes]\label{thm:sel}
Sous l'Hypothèse~\ref{hyp:H}, supposons $\Delta\ne0$ et $\bar\Delta\ne0$.
Notons $\gamma=\max(\rho,1/\rho)>1$, $\bar\gamma=\max(\bar\rho,1/\bar\rho)>1$.
Alors exactement un des trois régimes suivants a lieu :
\begin{itemize}
  \item $\Delta>0$ (donc $\bar\Delta<0$) $\iff aA+dD>bC+cB$ :
  $\lim_{N\to\infty}\limsup_{t\to\infty}\mathbb P((v_t,\vbar_t)\ne(\Inc,\Exc)) = 0$.
  \item $\bar\Delta>0$ (donc $\Delta<0$) $\iff dC+aB>cA+bD$ :
  $\lim_{N\to\infty}\limsup_{t\to\infty}\mathbb P((v_t,\vbar_t)\ne(\Exc,\Inc)) = 0$.
  \item $\Delta<0$ et $\bar\Delta<0$ :
  $\lim_{N\to\infty}\limsup_{t\to\infty}\mathbb P((v_t,\vbar_t)\ne(\Exc,\Exc)) = 0$.
\end{itemize}
\end{theorem}

\begin{proof}
Les équivalences découlent de la Proposition~\ref{prop:pv}. L'exclusion
mutuelle des trois régimes découle du Corollaire~\ref{cor:excl}. Pour la
borne : par le Théorème~\ref{thm:stat}, $\mathbb P(v_t\ne\Inc)\to\pi(\Exc)\le\gamma^{-N}$
si $\Delta>0$ (et symétriquement pour $\vbar_t$), sans invoquer de loi
stationnaire jointe pour la paire $(v_t,\vbar_t)$ -- celle-ci n'est pas
nécessairement irréductible. En prenant d'abord $\limsup_{t\to\infty}$ puis
l'union bound,
\[
  \limsup_{t\to\infty}\mathbb P((v_t,\vbar_t)\ne(\Inc,\Exc))
  \le \limsup_{t\to\infty}\mathbb P(v_t\ne\Inc)+\limsup_{t\to\infty}\mathbb P(\vbar_t\ne\Exc)
  \le \gamma^{-N}+\bar\gamma^{-N},
\]
sans hypothèse d'indépendance entre $v_t$ et $\vbar_t$. Puis
$\gamma^{-N}+\bar\gamma^{-N}\to0$ quand $N\to\infty$. Les deux autres
régimes se traitent de façon identique.
\end{proof}

\section{Concentration structurelle multi-features}

\begin{theorem}\label{thm:multi}
Clause à $n$ features gouvernées par $2n$ automates, en préférence
stationnaire non dégénérée pour chacun ($r_j=p_j^+/p_j^-\ne1$),
$\gamma_j=\max(r_j,1/r_j)>1$, $\gamma_{\min}=\min_j\gamma_j$. Alors, sans
hypothèse d'indépendance entre les automates,
\[
  \limsup_{t\to\infty}\mathbb P(C_t\ne C^\star) \le \sum_{j=1}^{2n}\frac1{1+\gamma_j^N}
  \le 2n\,\gamma_{\min}^{-N} \xrightarrow[N\to\infty]{} 0,
\]
où $C^\star$ désigne la configuration de littéraux stationnairement
préférée par la dynamique.
\end{theorem}

\begin{proof}
Soit $A_j$ l'évènement « l'automate $j$ est du mauvais côté ». Si tous les
$A_j$ sont faux, $C_t=C^\star$ ; par contraposée,
$\{C_t\ne C^\star\}\subseteq A_1\cup\cdots\cup A_{2n}$, d'où par l'union
bound : $\mathbb P(C_t\ne C^\star)\le\sum_j\mathbb P(A_j)$. Par le
Théorème~\ref{thm:stat}, $\mathbb P(A_j)=1/(1+\gamma_j^N)$ (que $r_j>1$ ou
$r_j<1$), d'où la première inégalité. Ensuite $1/(1+\gamma_j^N)<\gamma_j^{-N}\le\gamma_{\min}^{-N}$
pour chacun des $2n$ termes.
\end{proof}

\section{Identification sous dépendance de features arbitraire}

\begin{definition}\label{def:ctrue}
Soit $C_{\mathrm{true}}(x)=\bigwedge_{j\in J_+}x_j\wedge\bigwedge_{j\in J_-}(1-x_j)$
une conjonction fixe de $m=|J_+|+|J_-|\ge1$ littéraux, $J_+\cap J_-=\varnothing$.
On pose $Y^\star=C_{\mathrm{true}}(X)$.
\end{definition}

\begin{theorem}\label{thm:identif}
Sous l'Hypothèse~\ref{hyp:H}, supposons $Y=Y^\star$ p.s. -- la loi jointe de
$(X_1,\ldots,X_n)$ est quelconque (dépendante, déséquilibrée). Pour chaque
variable $j$, notons $A_j=P(X_j{=}0,Y{=}0)$, $B_j=P(X_j{=}1,Y{=}0)$,
$C_j=P(X_j{=}0,Y{=}1)$, $D_j=P(X_j{=}1,Y{=}1)$. Alors :
\begin{itemize}
  \item pour $j\in J_+$ : si $aA_j+dD_j>cB_j$, alors $\Delta_j>0$ et $\bar\Delta_j<0$ ;
  \item pour $j\in J_-$ : si $aB_j+dC_j>cA_j$, alors $\bar\Delta_j>0$ et $\Delta_j<0$ ;
  \item pour $k\notin J_+\cup J_-$ : si $aA_k+dD_k<bC_k+cB_k$ \emph{et}
  $aB_k+dC_k<cA_k+bD_k$, alors $\Delta_k<0$ et $\bar\Delta_k<0$.
\end{itemize}
Si ces conditions sont vérifiées pour tous les littéraux concernés, alors
$C^\star=C_{\mathrm{true}}$.
\end{theorem}

\begin{proof}
Pour $j\in J_+$ : $Y{=}1\Rightarrow X_j{=}1$ (fait logique, indépendant de la
loi de $X$), donc $C_j=P(X_j{=}0,Y{=}1)=0$. Par la Proposition~\ref{prop:pv} :
$\Delta_j=aA_j+dD_j-bC_j-cB_j=aA_j+dD_j-cB_j$, positif ssi $aA_j+dD_j>cB_j$.
Par le Corollaire~\ref{cor:excl}, $\Delta_j>0\Rightarrow\bar\Delta_j<0$.

Pour $j\in J_-$ : symétriquement, $Y{=}1\Rightarrow X_j{=}0$ donne
$D_j=P(X_j{=}1,Y{=}1)=0$, d'où $\bar\Delta_j=dC_j+aB_j-cA_j-bD_j=aB_j+dC_j-cA_j$,
positif ssi $aB_j+dC_j>cA_j$, et alors $\Delta_j<0$ par le Corollaire~\ref{cor:excl}.

Pour $k$ non pertinent : $A_k,B_k,C_k,D_k$ sont quelconques, aucune
simplification logique n'est possible, donc on requiert directement
$\Delta_k=aA_k+dD_k-bC_k-cB_k<0$ et $\bar\Delta_k=dC_k+aB_k-cA_k-bD_k<0$.
\end{proof}

\section{Robustesse à un bruit de label}

\begin{proposition}[Réécriture comme espérance]\label{prop:c-func}
Avec $c,\bar c:\{0,1\}^2\to\mathbb R$ définies par
\[
  c(0,0)=a,\ c(0,1)=-b,\ c(1,0)=-c,\ c(1,1)=d, \qquad
  \bar c(0,0)=-c,\ \bar c(0,1)=d,\ \bar c(1,0)=a,\ \bar c(1,1)=-b,
\]
on a $\Delta_j=\mathbb E[c(X_j,Y)]$ et $\bar\Delta_j=\mathbb E[\bar c(X_j,Y)]$
pour toute loi $P(X,Y)$.
\end{proposition}

\begin{proof}
Identification directe : $c(x,y)$ est la valeur signée de l'action sur
$v_j$ dans la Définition~\ref{def:regle} au cas $(x,y)$, et
$\mathbb E[c(X_j,Y)]=\sum_{x,y}c(x,y)P(X_j{=}x,Y{=}y)$ redonne exactement
$aA_j+dD_j-bC_j-cB_j=\Delta_j$ (Proposition~\ref{prop:pv}). Même calcul pour $\bar c$.
\end{proof}

\begin{lemma}[Saut ponctuel maximal]\label{lem:jump}
Pour $X_j$ fixé et $Y\ne Y'$ dans $\{0,1\}$ :
\[
  |c(X_j,Y)-c(X_j,Y')|\le M, \qquad |\bar c(X_j,Y)-\bar c(X_j,Y')|\le M,
  \qquad M:=\max(a+b,\ c+d).
\]
\end{lemma}

\begin{proof}
$|c(0,0)-c(0,1)|=|a-(-b)|=a+b$ ; $|c(1,0)-c(1,1)|=|-c-d|=c+d$. Le maximum
sur les deux valeurs de $X_j$ est $M$. Calcul identique pour $\bar c$ (mêmes
deux valeurs, échangées entre $X_j=0$ et $X_j=1$).
\end{proof}

\begin{theorem}[Robustesse au bruit]\label{thm:bruit}
Sous l'Hypothèse~\ref{hyp:H}, reprenons les notations du
Théorème~\ref{thm:identif} et notons
$\Delta_j^\star,\bar\Delta_j^\star,\Delta_k^\star,\bar\Delta_k^\star$ les
dérives calculées avec le label sans bruit $Y^\star$. Soit
\[
  \mu = \min\Bigl(
    \{\Delta_j^\star : j\in J_+\} \cup \{-\bar\Delta_j^\star : j\in J_+\} \cup
    \{-\Delta_j^\star : j\in J_-\} \cup \{\bar\Delta_j^\star : j\in J_-\} \cup{}
  \]
  \[
    \{-\Delta_k^\star : k\notin J_+\cup J_-\} \cup \{-\bar\Delta_k^\star : k\notin J_+\cup J_-\}
  \Bigr),
\]
et supposons $\mu>0$ (les conditions du Théorème~\ref{thm:identif} sont
satisfaites sans bruit). Soit $Y$ un label tel que $\eta:=P(Y\ne Y^\star)>0$,
sans aucune autre hypothèse sur la loi du bruit. Si
\[
  \boxed{\eta < \frac{\mu}{M}, \qquad M=\max(a+b,\,c+d),}
\]
alors toutes les inégalités du Théorème~\ref{thm:identif} restent
satisfaites avec $Y$ à la place de $Y^\star$, et $C^\star=C_{\mathrm{true}}$.
\end{theorem}

\begin{proof}
\emph{Étape 1 (espérance).} $\Delta_j=\mathbb E[c(X_j,Y)]$ par la
Proposition~\ref{prop:c-func} ; notons $\Delta_j^\star=\mathbb E[c(X_j,Y^\star)]$.

\emph{Étape 2 (couplage).} $X$ suit la même loi dans les deux cas (seul le
label diffère) ; couplons $Y$ et $Y^\star$ sur le même espace de
probabilité que $X$, de sorte que $\{Y\ne Y^\star\}$ soit exactement
l'évènement de bruit, de probabilité $\eta$. Alors
\[
  \Delta_j - \Delta_j^\star = \mathbb E\bigl[(c(X_j,Y)-c(X_j,Y^\star))\mathbf 1[Y\ne Y^\star]\bigr],
\]
le terme entre crochets étant nul dès que $Y=Y^\star$.

\emph{Étape 3 (saut ponctuel).} Par le Lemme~\ref{lem:jump},
$|c(X_j,Y)-c(X_j,Y^\star)|\le M$ sur $\{Y\ne Y^\star\}$, pour toute valeur de $X_j$.

\emph{Étape 4 (passage à l'espérance).} En combinant les étapes 2 et 3,
\[
  |\Delta_j-\Delta_j^\star| \le M\,\mathbb E[\mathbf 1[Y\ne Y^\star]] = M\eta,
\]
et de même $|\bar\Delta_j-\bar\Delta_j^\star|\le M\eta$, $|\Delta_k-\Delta_k^\star|\le M\eta$,
$|\bar\Delta_k-\bar\Delta_k^\star|\le M\eta$ (mêmes fonctions $c,\bar c$).

\emph{Étape 5 (préservation des signes).} Soit $\eta<\mu/M$, i.e.\ $M\eta<\mu$.
Pour $j\in J_+$ : $\Delta_j^\star\ge\mu$ par définition de $\mu$, donc
$\Delta_j\ge\Delta_j^\star-M\eta\ge\mu-M\eta>0$. De même
$-\bar\Delta_j^\star\ge\mu$ donne $\bar\Delta_j\le\bar\Delta_j^\star+M\eta\le-\mu+M\eta<0$.
Le même calcul, terme à terme, s'applique à $j\in J_-$ et aux $k$ non
pertinents. Toutes les inégalités strictes du Théorème~\ref{thm:identif}
restent donc vraies avec $Y$ à la place de $Y^\star$.
\end{proof}

\end{document}
